\documentclass[10pt,a4paper]{amsart}
\usepackage[margin=19mm]{geometry}
\usepackage{amsmath,amssymb,mathtools}
\usepackage[scr=rsfs,cal=euler]{mathalfa}
\usepackage{xfrac}
\usepackage[unicode,hidelinks,bookmarks=false]{hyperref}
\newcommand{\ie}{i.e.\ }
\newcommand{\cf}{cf.\ }
\newcommand{\resp}{respectively\ }
\newcommand{\Subsection}{\subsection}
\providecommand{\nobreakdash}{\nobreak\hbox{-}}
\setlength{\emergencystretch}{2em}
\multlinegap0pt
\usepackage{multirow}
\usepackage{amssymb}
\usepackage{xcolor}
\usepackage{booktabs}
\usepackage[normalem]{ulem}
\usepackage{tikz}
\newtheorem{proposition}{Proposition}
\newtheorem{definition}{Definition}
\newtheorem{theorem}{Theorem}
\newcommand{\N}{\mathbb{N}} % symbol for natural numbers
\newcommand{\cG}{{\mathcal G}}         %% wird umdefiniert s.u.
\def\none{\tilde{\mbox{$1 \hspace{-1.0mm}  {\bf l}$}}}
\newcommand{\quadtext}[1]{\quad\text{#1}\quad}
\title{Asymptotic robustness of entanglement in noisy quantum networks and graph connectivity}
\author{Fernando Lled\'o, Carlos Palazuelos, Julio I. de Vicente}
\date{Source: arXiv:2411.12548v2 (CC BY 4.0)}
\begin{document}
\maketitle

\noindent\emph{Source and licence.} Asymptotic robustness of entanglement in noisy quantum networks and graph connectivity. Version arXiv:2411.12548v2 (CC BY 4.0), distributed by arXiv under the Creative Commons Attribution 4.0 International licence, http://creativecommons.org/licenses/by/4.0/, as recorded in the arXiv licence metadata for this exact version and shown on the abstract page https://arxiv.org/abs/2411.12548v2. Full licence notice: \texttt{licenses/arxiv-2411.12548-CC-BY-4.0.txt}.\par\smallskip

\noindent\textit{Editorial scope.} Counted unit: the article's theorem th:degABS (source0.tex lines 1044--1047). The article's complete original proof of it is delivered with it (source0.tex lines 1049--1175, one \texttt{proof} environment of 127 lines). No mathematics is written, repaired or re-derived: the statement and the proof are the article's own lines, in the article's own notation, and no retained line is substituted or re-flowed.  Retained uncounted as the source-local context the proof needs: lines 277--277; lines 293--295; lines 308--310; lines 446--448; lines 503--511; lines 572--581.  Nothing was omitted from any retained span, so the package carries no omission record.  No retained line contains a citation, so the wrapper carries no bibliography and no bibliography entry.  No retained line contains a figure environment, an image-inclusion command or drawing code, so no image is delivered and the package declares no asset dependency.  Editorial formatting only, in addition: an AMS \texttt{amsart} preamble that restates the article's own declarations and macros used by the retained lines, namely the environments \texttt{proposition}, \texttt{definition}, \texttt{theorem} on separate counters, together with the standard packages those lines need.  The retained environments print in this document as Proposition 1, Definition 1, Theorem 1 in order of appearance, whereas the article prints the same items with the numbering of its own full document.  The complete native source of the article as distributed in the arXiv e-print for this exact version is bundled unchanged as \texttt{sources/168762/source0.tex}.
\subsection{Multipartite quantum systems}\label{sec21}

\begin{equation}\label{isotropicd}
\rho(p,d)=p\phi^+_d+(1-p)\none_{d^2},
\end{equation}

\begin{equation}\label{isotropicthreshold}
p>\frac{1}{d+1}.
\end{equation}

\begin{equation}\label{networkGdp}
\sigma_d(G,p)= \bigotimes_{e\in E}\rho_e(p,d),
\end{equation}

\begin{proposition}\label{prop:facts}
Let $G$ be a graph and consider a fixed dimension $d$. Then
\begin{itemize}
 \item[(i)] \label{fact1}
 there exists $\epsilon>0$ such that $\sigma_d(G,p)$ is GME if $p>1-\epsilon$;
 \item[(ii)] \label{fact2}
 if $\sigma_d(G,p)$ is BS, then $\sigma_d(G,p')$ is BS for all $p'\leq p$.
\end{itemize}
\end{proposition}

\begin{definition}\label{def:asymptotic}
 Consider a graph sequence $\cG=(G_n)_{n\in\N}$.
\begin{itemize}
 \item[(i)]%\label{def:AGME}
 $\cG$ is \textbf{asymptotically genuine multipartite entangled (AGME)} in a given dimension $d$ if there exists $p_0\in(1/(d+1),1)$ such that for all $p\geq p_0$ and for all $n\in\N$ the isotropic network state $\sigma_d(G_n,p)$ is GME. If the graph sequence $\cG$ is AGME for any given $d\geq2$, we say simply that $\cG$ is AGME.

 \item[(ii)]%\label{def:ABS}
$\cG$ is \textbf{asymptotically biseparable (ABS)} in a given dimension $d$ if for all $p\in(1/(d+1),1)$ there exists $n_0\in\mathbb{N}$ such that for all $n\geq n_0$ the state $\sigma_d(G_n,p)$ is BS. If $\cG$ is ABS for any given $d\geq2$, we say that $\cG$ is ABS.
\end{itemize}
\end{definition}

\begin{theorem}\label{th:degABS}
Any graph sequence $\cG=(G_n)_{n\in\N}$ for which the associated sequence of maximal degrees satisfies
$\delta_{\mathrm{max}}(G_n)\in o\left(\log|G_n|\right)$ is ABS.
\end{theorem}

\begin{proof}
For any given $d$, which is considered fixed in the following, we will simplify the notation of the isotropic
states in Eq.~(\ref{isotropicd}) putting $\phi^+$ and $\none$ instead of $\phi^+_d$ and $\none_{d^2}$, respectively, and use edge subindices to label the Hilbert spaces on which the operators act.

According to Definition~\ref{def:asymptotic}~(ii) we need to show that the isotropic network state
$\sigma_d(G_n,p)$ given in Eq.~(\ref{networkGdp}) is biseparable for all $p<1$ and for $n$ large enough.
Although it is not strictly necessary, we will show in order to simplify the reasoning that
$\sigma_d(G_n,p)$ is actually 1-biseparable. For this, rewrite first (using distributivity) the isotropic network state in terms of $2^{|E_n|}$ summands labeled by subsets of the edge set $E_n$ as
\begin{equation}\label{eq:summands}
\sigma_d(G_n,p)=\sum_{F_n\subseteq
E_n}p^{|F_n|}(1-p)^{|E_n|-|F_n|}\bigotimes_{e\in F_n}\phi^+_e\bigotimes_{e\notin
F_n}\none_e\;.
\end{equation}
We will split the preceding sum into three parts labeled by families of edge subsets
corresponding to different configurations of connectivity of vertices.
For this denote for any subset $F_n\subset E_n$
the edges in $F_n$ adjacent with a vertex $v\in V_n$ by
\begin{equation}
 F_n(v):=F_n\cap E_n(v)\;.
\end{equation}


\begin{itemize}
 \item We consider first all subsets of edges that label in the previous sum
\textit{non} 1-biseparable states:
\begin{equation}\label{bn}
 B_n=\{F_n\subseteq E_n \mid
 F_n(v)\neq\emptyset\quadtext{for any} v\in V_n\}.
\end{equation}

\item Next, we consider all subsets of edges labeling states that
are 1-biseparable \textit{exactly} at one vertex $v$:
\begin{equation}\label{an}
A_n=\{F_n'\subseteq E_n \mid \textrm{ there exists a unique }
      v\in V_n \textrm{ such that } F_n'(v)=\emptyset\}\;.
\end{equation}
In fact, if $F_n'\in A_n$ and for (exactly) one vertex $v\in V_n$
we have that $F_n'(v)=\emptyset$, then (see Subsection~\ref{sec21})
\begin{equation}
\bigotimes_{e\in {F'_n}}
\phi^+_e\bigotimes_{e\notin {F'_n}}\none_e\in S_v(H).
\end{equation}
Thus, mixtures of these terms for different
choices of $v$ give rise to
a 1-biseparable state.

\item Finally, we consider all remaining subsets of edges,
\begin{equation}\label{cn}
C_n= E_n\setminus \left(A_n\cup B_n\right)\;,
\end{equation}
which manifestly correspond to 1-biseparable terms.
\end{itemize}

Note that for any $F_n\in B_n$ and for any $v\in V_n$ there exists a
unique $F'_n\in A_n$ such that $F'_n(v)=\emptyset$ and $F_n(w)=F'_n(w)$ for all $w\neq v$ ($F'_n$ is obtained from $F_n$ by removing all edges adjacent to $v$).
The idea of the proof consists in recombining non-1-biseparable terms in $B_n$
with the 1-biseparable terms from $A_n$ (with appropriate weights) so
as to overall obtain a 1-biseparable state for $n$ sufficiently large. In this
process the summands labeled by $C_n$ will remain untouched as they already correspond to
1-biseparable states.
To establish the weights needed in the mixture notice that any $F_n\in B_n$ can be
combined with $|G_n|$ terms in $A_n$ with the preceding property (one for each choice of vertex). However, if we want to do so
for all terms in $B_n$, each term $F'_n\in A_n$ such that $F'_n(v)=\emptyset$
for a given $v\in V_n$ needs to be used for all terms $F_n\in B_n$ such that
$F_n(w)=F'_n(w)$ for all $w\neq v$ and note that there are $2^{\delta_n(v)}-1$
such terms (one for each possible choice of $F_n(v)$ excluding the empty
set).

Considering the vertex refined decomposition of the $B_n$ summands as
\begin{eqnarray}
\lefteqn{\sum_{F_n\in B_n}
p^{|F_n|}(1-p)^{|E_n|-|F_n|}\bigotimes_{e\in
F_n}\phi^+_e\bigotimes_{e\notin F_n}\none_e} \qquad\qquad\nonumber\\[3mm]
&=&
\sum_{F_n\in B_n}\sum_{v\in V_n}\frac{1}{|G_n|}\;\;
p^{|F_n|}(1-p)^{|E_n|-|F_n|}\bigotimes_{e\in
F_n}\phi^+_e\bigotimes_{e\notin F_n}\none_e\;
\end{eqnarray}
(and simmilarly for the $A_n$ summands with weights
$\left(2^{\delta_n(v)}-1\right)^{-1}$) we can rewrite Eq.~(\ref{eq:summands}) as
\begin{eqnarray}
\sigma_d(G_n,p)&=&
\sum_{F_n\in B_n}
p^{|F_n|}(1-p)^{|E_n|-|F_n|}\bigotimes_{e\in
F_n}\phi^+_e\bigotimes_{e\notin F_n}\none_e
+ \sum_{F_n'\in A_n}\Big(\cdots \Big)
+ \sum_{F_n\in C_n}\Big(\cdots \Big)
\nonumber \\[3mm]
&=& \label{regroup}
\sum_{F_n\in B_n}\sum_{v\in V_n}
p^{|F_n|}(1-p)^{|E_n|-|F_n|}
\chi_{(v)}^{(F_n)}
\bigotimes_{e\in F_n\setminus F_n(v)}\phi^+_e
\bigotimes_{e\notin F_n}\none_e \nonumber\\[2mm]
&& +\sum_{F_n\in
C_n}p^{|F_n|}(1-p)^{|E_n|-|F_n|}\bigotimes_{e\in
F_n}\phi^+_e\bigotimes_{e\notin F_n}\none_e,\label{eq:summands2}
\end{eqnarray}
where
\begin{equation}
\chi_{(v)}^{(F_n)}=
\frac{1}{|G_n|}\bigotimes_{e\in F_n(v)}\phi^+_e+\frac{p^{-|F_n(v)|}(1-p)^{|F_n(v)|}}{2^{\delta_n(v)}-1}\bigotimes_{e\in F_n(v)}\none_e
\end{equation}
is an unnormalized state for each choice of $F_n\in B_n$ and $v\in V_n$ realizing the combination of $B_n$ and $A_n$ summands at each vertex. In fact, up to proper normalization, $\chi_{(v)}^{(F_n)}$ is an isotropic state in $H_1\otimes H_2$ where $\dim H_1=\dim H_2=d^{F_n(v)}$.

Now, according to Eq.~(\ref{isotropicthreshold}), such state is separable if and only if
\begin{equation}
\frac{\frac{1}{|G_n|}}{\frac{1}{|G_n|}+\frac{p^{-|F_n(v)|}(1-p)^{|F_n(v)|}}{2^{\delta_n(v)}-1}}\leq \frac{1}{d^{|F_n(v)|}+1}
\end{equation}
or, equivalently,
\begin{equation}
d^{|F_n(v)|}\Big(\frac{p}{1-p}\Big)^{|F_n(v)|}(2^{\delta_n(v)}-1)\leq |G_n|.
\end{equation}

We want to show that for every $p\in (0,1)$, there exists $n_0\geq 1$ such that for every $n\geq n_0$ the previous state is separable for all $F_n$ and $v$. Actually, according to Proposition~\ref{prop:facts}, we can assume that $p>1/2$, so that $p/(1-p)>1$.

Hence, it suffices to show that  for every $p\in (1/2,1)$, there exists $n_0\geq 1$ such that for every $n\geq n_0$ we have
\begin{equation}
d^{\delta_\mathrm{max}(G_n)}\Big(\frac{p}{1-p}\Big)^{\delta_\mathrm{max}(G_n)}2^{\delta_\mathrm{max}(G_n)}\leq |G_n|,
\end{equation}
or, equivalently,
\begin{equation}
\Big(\log d+ \log \frac{p}{1-p}+\log 2\Big)
  \delta_\mathrm{max}(G_n)\leq \log |G_n|,
\end{equation}
which follows from our assumption on the growth of $\delta_\mathrm{max}(G_n)$.
\end{proof}

\end{document}
